\subsection{مبانی نظری و فرمول‌بندی ریاضی}
فرض کنید بردارهای پرس‌وجو و کلید در یک فضای $d$-بعدی به صورت $\mathbf{q}, \mathbf{k} \in \mathbb{R}^d$ قرار داشته باشند. در روش درهم‌سازی مبتنی بر ابرصفحه‌های تصادفی زاویه‌ای، به ازای هر دور $r \in \{1, \dots, R\}$، یک ماتریس تصویر تصادفی $\mathbf{R}_r \in \mathbb{R}^{d \times \frac{b}{2}}$ با توزیع نرمال استاندارد نمونه‌برداری و بهنجار می‌شود:
\begin{equation}
\tilde{\mathbf{R}}_r = \frac{\mathbf{R}_r}{\|\mathbf{R}_r\|_2}
\end{equation}
سپس، شناسه سطل هر بردار با درهم‌سازی متقارن تصادفی تعیین می‌گردد:
\begin{equation}
h_r(\mathbf{x}) = \operatorname{argmax}\left( [\mathbf{x} \tilde{\mathbf{R}}_r ; -\mathbf{x} \tilde{\mathbf{R}}_r] \right) \in \{0, \dots, b - 1\}
\end{equation}
که در آن $b$ تعداد سطل‌های درهم‌سازی است. احتمال قرارگیری دو بردار در یک سطل رابطه‌ای مستقیم با زاویه کسینوسی $\theta$ میان آنها دارد:
\begin{equation}
\Pr[h(\mathbf{q}) = h(\mathbf{k})] = 1 - \frac{\theta(\mathbf{q}, \mathbf{k})}{\pi}
\end{equation}

برای افزایش ضریب اطمینان و جلوگیری از حذف تصادفی توکن‌های مرتبط، الگوریتم به تعداد $R$ دور مستقل اجرا می‌شود. ماتریس امتیازات درون هر سطل به شکل زیر تعریف می‌گردد:
\begin{equation}
S_{i, j}^{(r)} = \begin{cases}
\frac{\mathbf{q}_i \mathbf{k}_j^T}{\sqrt{d}} & \text{if } h_r(\mathbf{q}_i) = h_r(\mathbf{k}_j) \text{ and } m_j = 1 \\
-\infty & \text{otherwise}
\end{cases}
\end{equation}
خروجی نهایی از میانگین‌گیری نتایج تمامی دورهای درهم‌سازی حاصل می‌شود:
\begin{equation}
\mathbf{O} = \frac{1}{R} \sum_{r=1}^R \operatorname{Softmax}\left(\mathbf{S}^{(r)}\right) \mathbf{V}
\end{equation}

که در آن $\sigma_i(\mathbf{W})$ بیانگر $i$-امین مقدار منفرد مرتب‌شده ماتریس به صورت نزولی است. نسبت حفظ انرژی طیفی $\eta_r = \frac{\sum_{i=1}^r \sigma_i^2}{\sum_{i=1}^{d} \sigma_i^2}$ تضمین می‌کند که با انتخاب آستانه $\tau \ge 0.90$، بیش از نود درصد از پراکندگی واریانس داده‌ها حفظ شده و افت سیگنال در انتشار گرادیان رو به عقب به حداقل ممکن کاهش می‌یابد.

\paragraph{تحلیل مجانبی و پیچیدگی اشتراک حافظه سخت‌افزاری}
از منظر پیچیدگی محاسباتی و انتقال داده‌ها در سلسومراتب حافظه (حافظه سراسری \lr{HBM} به حافظه محلی \lr{SRAM})، عملیات ضرب ماتریسی در این مکانیزم نیازمند تخصیص بافر موقت با ابعاد فاکتورگیری‌شده است. برای پردازش دنباله‌ای به طول $N$ با بعد مخفی $D$، بار محاسباتی هسته از فرم درجه دو به رابطه زیر تقلیل می‌یابد:
\begin{equation}
\label{eq:complexity_lsh_attention}
\mathcal{C}_{\text{lsh\_attention}} = \mathcal{O}\left( B \cdot N \cdot D \cdot r + B \cdot H \cdot N \cdot d_k \right) \ll \mathcal{O}\left( B \cdot H \cdot N^2 \cdot d_k \right)
\end{equation}
این بهبود مرتبه پیچیدگی، با کاهش چشمگیر حجم تبادل تانسورها در گذرگاه‌های ارتباطی کارت گرافیک، مانع از اشباع پهنای باند حافظه (پدیده \lr{Memory-Bound}) در فرآیند استنتاج طولانی‌مدت می‌گردد.
